%% kapa: exam=1 round=34 year=2023 period=1 form=A subject=economics subject_ko=경제학원론 seq=34 number=74
%% source: Q-Net 감정평가사 기출문제 (cst00309) artlSeq=5211789 file=raw/1008/5211789/2023년도 제34회 감정평가사 1차시험 1교시.pdf
%% conversion: kordoc 4.24.51 -> tools/gen1.py (auto); review=proofread-sonnet; flags=none
\begin{question}{74}
필립스곡선에 관한 설명으로 옳지 않은 것은?
\begin{choices}
  \item 필립스(A. W. Phillips)는 임금상승률과 실업률간 음(-)의 경험적 관계를 발견했다.
  \item 우상향하는 단기 총공급곡선과 오쿤의 법칙(Okun's Law)을 결합하면 필립스곡선의 이론적 근거를 찾을 수 있다.
  \item 적응적 기대를 가정하면 장기에서도 필립스곡선은 우하향한다.
  \item 단기 총공급곡선이 가파른 기울기를 가질수록 필립스곡선은 가파른 기울기를 가진다.
  \item 새고전학파(New Classical)는 합리적 기대를 가정할 경우 국민소득의 감소 없이 인플레이션을 낮출 수 있다고 주장한다.
\end{choices}
\end{question}
\answer{3}
